\section{The Sin and Punishment of the Fallen Angels}\label{sec:st-fall}

How can a creature of pure intelligence---with no inflamed passions to cloud its sight, no prior bad habit to bend its will, and no possibility of erring about natural truth---ever commit a mortal sin? When human beings sin, disordered flesh or partial ignorance ordinarily darkens practical judgment; an unfallen angel had neither. In Questions~63 and~64 of the \emph{Prima Pars}, deepened a few years later in his disputed questions \emph{De malo} (q.~16, aa.~2--6), St.~Thomas Aquinas resolves the paradox with severe clarity: the angel did not mistake evil for good, nor did its intellect fail before its will fell; it chose a genuine spiritual good while freely declining to attend to the higher rule of God's will.

\subsection{The Possibility, Nature, and Object of Angelic Sin (\ST{I, q.~63, aa.~1--4})}

Sacred Scripture testifies directly to angelic guilt: ``in his angels he found wickedness'' (Job~4:18). Aquinas grounds the possibility of angelic sin in the very metaphysics of creaturehood (\ST{I, q.~63, a.~1}). Whether in nature, art, or morals, sin is deviation from the due rule of action. ``Were the craftsman's hand the rule itself engraving, he could not engrave the wood otherwise than rightly; but if the rightness of engraving be judged by another rule, then the engraving may be right or faulty.'' Only the Divine will is its own rule, because God is ordered to no higher end. Every created will, angelic or human, has rectitude solely insofar as it regulates itself by the Divine will; impeccability therefore belongs to God alone by nature, and to any creature solely as a gift of grace.

In \ST{I, q.~63, a.~1, ad~4}---expanded in \emph{De malo}~(q.~16, a.~2, ad~7)---Aquinas solves the puzzle of sin without passion or prior error. Mortal sin happens in two ways. Often a human being chooses something evil in itself (such as adultery), which always requires a particular practical error induced by passion or habit. That road was closed to the unfallen angel. Yet sin also arises when a free will chooses something genuinely good in itself while omitting the proper rule and measure---as if a person were to pray while refusing the order appointed by the Church. Such a choice presupposes no prior ignorance or intellectual error, solely a free defect of negation: ``absence of consideration of the things which ought to be considered,'' failing to govern one's choice by the higher rule of Divine Wisdom (\emph{De malo}~q.~16, a.~2).

What sins can a bodiless spirit actually commit (\ST{I, q.~63, a.~2})? Aquinas distinguishes between guilt (\emph{reatus}) and personal affection (\emph{affectus}). By guilt, demons incur every kind of sin because they tempt men to commit every sin. By personal affection, however, only two spiritual sins can exist in an angel: \emph{pride} first, and \emph{envy} as its consequence. St.~Augustine puts the matter bluntly in \emph{City of God}~(XIV.3): the devil ``is not a fornicator nor a drunkard, nor anything of the like sort; yet he is proud and envious.'' Because pure spirits cannot be drawn by fleshly pleasures---the ground on which Augustine definitively set aside the early Watchers hypothesis of carnal angelic sin (\emph{City of God}~XV.23, contrasting St.~Justin Martyr, \emph{2~Apology}~5)---a spirit can sin in desiring spiritual good only by refusing due subjection to a superior. That refusal is pride. From coveting singular excellence follows envy: grieving over another's good as a rival to one's own pre-eminence, including grieving over human salvation and over God's glory in raising man up (\ST{I, q.~63, a.~2, ad~1--3}).

Isaias captures the inner voice of that pride: ``I will ascend into heaven\dots{} I will be like the most High'' (Is~14:13--14). Yet what did Lucifer actually desire (\ST{I, q.~63, a.~3})? Aquinas draws a razor-sharp distinction between \emph{equality} with God and \emph{likeness} to God. Lucifer could not desire absolute equality with God: by natural self-knowledge he saw clearly that equality with the Creator is impossible, and even were it possible, passing into a higher species would annihilate his own being---just as an ass does not desire to become a horse, because every creature naturally wills to preserve its own nature (\ST{I, q.~63, a.~3}; \emph{De malo}~q.~16, a.~3, distinguishing true choice from mere conditional wishing, \emph{velleitas}, in ad~9).

Instead, Lucifer sinned by coveting an inordinate \emph{likeness} to God in the supernatural order: seeking final beatitude from his own natural power rather than as a gift of divine grace. Whether he turned away from supernatural beatitude to rest in natural perfection alone, or coveted the supernatural vision of God through his own strength---``by the power of his own nature; not indeed without God operating in nature, yet without God conferring grace'' (\emph{De malo}~q.~16, a.~3)---both ways coincide with St.~Anselm's formula quoted by Aquinas (\emph{De casu diaboli}~IV): ``he sought that to which he would have come had he stood fast.'' Because what exists of itself naturally rules over others, coveting self-sufficient beatitude led Lucifer immediately to covet autonomous dominion over fellow angels (\ST{I, q.~63, a.~3}; \emph{De malo}~q.~16, a.~3, ad~4; cf.~St.~Bonaventure, \emph{Breviloquium}~II.7, noting that by presuming on his private good Lucifer made himself his own principle and end). Later scholastic schools explored two other lines: Bl.~John Duns Scotus (\emph{Ordinatio}~II, d.~6, q.~2) traced Lucifer's first sin to an immoderate self-love for his own beatitude (\emph{affectio commodi}, a kind of \emph{luxuria spiritualis}) unmoderated by justice (\emph{affectio iustitiae}), whereas Francisco Su\'arez (\emph{De angelis}~VII.13) and a later Scotist development centered Lucifer's revolt on proud refusal to adore the revealed mystery of the Incarnation and the hypostatic union of human nature with the Word.

No demon, however, is naturally wicked (\ST{I, q.~63, a.~4}). Against Manichaean and Cathar dualism, Pseudo-Dionysius affirms that ``the demons are not naturally wicked'' (\emph{De divinis nominibus}~IV), and the Fathers teach unanimously that evil is a privation introduced by free choice rather than a created essence (St.~Irenaeus, \emph{Adversus haereses}~IV.41.1; St.~Augustine, \emph{City of God}~XII.1, 6; St.~John of Damascus, \emph{De fide orthodoxa}~II.4). The Fourth Lateran Council solemnly defined the doctrine in \emph{Firmiter credimus} (DS~800): \emph{Diabolus enim et alii daemones a Deo quidem natura creati sunt boni, sed ipsi per se facti sunt mali}, ``For the devil and the other demons were indeed created by God naturally good, yet they became evil by themselves''---reaffirmed by the Council of Florence (\emph{Cantate Domino}, DS~1333, 1336). Whereas a beast's sensory appetite pursues particular goods linked to another creature's harm, an intellectual nature inclines to universal good, with which no evil can be joined; hence Porphyry erred in thinking demons naturally deceitful only because he imagined them as sensitive aerial animals (St.~Augustine, \emph{City of God}~X.11; \ST{I, q.~63, a.~4, ad~1--3}).

\subsection{The Instant, Leader, and Extent of the Fall (\ST{I, q.~63, aa.~5--9})}

Could the devil have sinned in the very first instant of his creation (\ST{I, q.~63, a.~5})? Scripture's verdict on the goodness of creation---``God saw all the things that he had made, and they were very good'' (Gen~1:31)---together with Isaias~14:12 and Ezechiel~28:13--15 (``Thou wast perfect in thy ways from the day of thy creation, until iniquity was found in thee''), led Aquinas in both the \emph{Summa} and \emph{De malo} to answer no. St.~Augustine had left the question undecided---inclining toward first-instant voluntary sin in \emph{De Genesi ad litteram}~(XI.16.21) while inclining against it in \emph{City of God}~(XI.13, 15), noting that voluntary first-instant sin is distinct from Manichaeism. When certain contemporary masters nevertheless asserted first-instant sin, ``this position was condemned by all the masters then lecturing at Paris'' (\emph{De malo}~q.~16, a.~4).

Between the \emph{Summa theologiae} and \emph{De malo}, however, Aquinas executes a striking doctrinal refinement in his proof. In \ST{I, q.~63, a.~5}, he argues from \emph{efficient causality}: although creation and an angelic act of will are both instantaneous, any operation occurring in the very first instant of a creature's existence flows from its producing agent---just as fire's upward motion comes from its generator, or a congenital limp from a defective generator. Because God cannot cause sin, the devil could not sin in instant~1. Yet because merit comes from God's grace, Aquinas concludes in \ST{I, q.~63, a.~5, ad~4} that all angels, created in grace, actually \emph{merited} in the first instant, until the rebels immediately placed an obstacle in the second instant and destroyed their preceding merit.

In his later \emph{De malo} (q.~16, a.~4), Aquinas candidly owns the difficulty of the proof (\emph{etsi difficile sit}, ``even though it be difficult'') and abandons the efficient-causality argument in favor of \emph{cognitive non-simultaneity}. Because an angel cannot actuate several distinct intelligible species in a single instant (\ST{I, q.~58, a.~2}), and supernatural realities of grace stand even farther from connatural knowledge than distinct natural objects stand from one another, the newly created angel could not simultaneously turn toward its own natural being and toward supernatural grace. Moving first to what is connatural, every angel in instant~1 turned to natural self-knowledge and natural love of God as Author of nature---acts in which no sin can arise---and only in the second discrete instant of angelic time turned toward or away from supernatural beatitude. Aquinas states the revised conclusion unequivocally: \emph{Et ideo Angelus in primo instanti suae creationis non fuit neque beatus per conversionem perfectam in Deum, neque peccator per aversionem ab ipso}, ``and therefore in the first instant of its creation, the angel was neither beatified through perfect conversion to God, nor a sinner through aversion from Him'' (\emph{De malo}~q.~16, a.~4). In the same article (\emph{De malo}~q.~16, a.~4, ad~16), he adds that mankind was not created merely to fill the void left by fallen spirits, for human beings would have been created ``even if the angelic ruin had never occurred'' (\emph{etiam si nunquam fuisset ruina angelica}).

Was there an extended interval (\emph{mora}) between creation and the fall (\ST{I, q.~63, a.~6})? While acknowledging that an interval is conceivable if angels were not created in grace---a view favored by St.~Bonaventure (\emph{Breviloquium}~II.7) and several Fathers---Aquinas judges it more probable that the fall happened immediately after the first instant, in the second discrete instant of angelic time. In instant~1, all angels turned to themselves in evening knowledge, which was good in all; in instant~2, the faithful turned in morning knowledge to the praise of the Word, while the rebels grew absorbed in themselves, ``swelling up with pride,'' and became night (\ST{I, q.~63, a.~6, ad~4}, citing St.~Augustine, \emph{De Genesi ad litteram}~IV.24; cf.~John~8:44, ``he stood not in the truth'').

Was the leader of the fallen angels the very highest of all angels (\ST{I, q.~63, a.~7})? Aquinas sets two great Fathers side by side with admirable balance. Looking to \emph{proneness} to sin, St.~John of Damascus held that the fallen chief had been entrusted with presiding over the earthly realm (\emph{De fide orthodoxa}~II.4)---an opinion Aquinas notes accords with God's governance of bodies through lower angels (St.~Augustine, \emph{De Trinitate}~III.4--5) and does not conflict with faith. Looking instead to the \emph{motive} for sin, however, Aquinas follows St.~Gregory the Great (\emph{Homiliae in Evangelia}~XXXIV; echoed in Peter Lombard, \emph{Sententiae}~II, d.~6): because angelic sin sprang not from frailty but from free pride sparked by excellence, it is more probable that Lucifer surpassed all other angels in natural brightness, without prejudicing Damascene's view. Ezechiel~28:14 calls him a ``covering cherub'' rather than a seraph because \emph{Cherubim} denotes fullness of knowledge, which can coexist with mortal sin, whereas \emph{Seraphim} denotes burning charity, which cannot (\ST{I, q.~63, a.~7, ad~1}).

Lucifer's sin caused the sin of the other fallen angels---not by compulsion, which no creature can impose upon a free will, but by persuasion and exhortation (\ST{I, q.~63, a.~8}). The Apocalypse depicts the dragon drawing ``the third part of the stars of heaven'' (Apoc~12:4), and Christ speaks of everlasting fire ``prepared for the devil and his angels'' (Matt~25:41), because divine justice subjects those who consent to evil inducement to the one who overcame them (2~Pet~2:19). Stripping away human delay for spoken words and deliberation, Lucifer manifested his proud will by intelligible speech and the consenting spirits agreed in the exact same instant (\ST{I, q.~63, a.~8, ad~1}, harmonizing Damascene's statement that all fell together). Proud subordinates accepted Lucifer's headship because pride willingly accepts a familiar natural superior over God when hoping to win beatitude by natural power (\ST{I, q.~63, a.~8, ad~2}); and acting with unretarded natural energy, the highest spirit fell into the deepest malice (\ST{I, q.~63, a.~8, ad~3}).

Finally, more angels stood firm than fell (\ST{I, q.~63, a.~9}). Reading Eliseus's reassurance---``Fear not: for there are more with us than with them'' (4~Kgs [2~Kgs]~6:16)---as referring to the holy angels outnumbering the demons, Aquinas reinforces the reading from nature itself: sin runs against natural inclination, and what opposes nature occurs less frequently than what accords with it. Whereas bodily passions draw the majority of men away from reason (\ST{I, q.~63, a.~9, ad~1--2}), pure spirits face no such drag. Under Gregory's view, some angels probably fell from every choir---though Scripture never calls demons ``Seraphim'' or ``Thrones,'' names of charity and divine indwelling, only ``Cherubim,'' ``Powers,'' and ``Principalities'' (\ST{I, q.~63, a.~9, ad~3}).

\angeldossier{\ST{I, q.~63, aa.~1--9} (Malice of the angels with regard to sin: possibility of sin, pride and envy, desiring likeness to God, non-natural malice, impossibility of first-instant sin, immediate fall, Lucifer's rank, causal persuasion, and proportion of the fallen).}{Defined (aa.~1, 4); Common teaching (aa.~2--3, 8); Thomist position (aa.~5--7, 9).}{Gen~1:31; 4~Kgs [2~Kgs]~6:16; Job~4:18; Ps~103:26 [Heb~104:26] (\emph{Glossa}); Eccles~1:15; Wis~12:10; Is~14:12--14; Ezech~28:13--15; Matt~25:41; John~8:44; 2~Pet~2:19; Apoc~12:4; Fourth Lateran Council, \emph{Firmiter credimus} (DS~800); Council of Florence, \emph{Cantate Domino} (DS~1333, 1336); St.~Irenaeus, \emph{Adversus haereses}~IV.40--41; Pseudo-Dionysius, \emph{De divinis nominibus}~IV and \emph{De caelesti hierarchia}~VI--VII; St.~Augustine, \emph{De civitate Dei}~VII, X.11, XI.13, 15, XII.1, 6, XIV.3, XV.23, \emph{De Genesi ad litteram}~I.15, II.8, IV.24, XI.16, and \emph{De Trinitate}~III.4--5; Ambrosiaster (\emph{De quaestionibus Veteris Testamenti}~CXIII); St.~Gregory the Great, \emph{Homiliae in Evangelia}~XXXIV and \emph{Moralia in Iob}~XXXI; St.~John of Damascus, \emph{De fide orthodoxa}~II.4; St.~Anselm of Canterbury, \emph{De casu diaboli}~IV (cited by Aquinas); Peter Lombard, \emph{Sententiae}~II, d.~6; Boethius, \emph{De hebdomadibus}; Aristotle, \emph{Physics}~VI, \emph{Metaphysics}~IX, and \emph{Nicomachean Ethics}~II.6; St.~Thomas Aquinas, \emph{De malo}~q.~16, aa.~2--4.}{Manichaean and Cathar dualism of an evil nature (aa.~1, 4); the early Watchers interpretation of carnal sin (St.~Justin Martyr, \emph{2~Apology}~5), definitively set aside by St.~Augustine (\emph{City of God}~XV.23) (a.~2); Bl.~John Duns Scotus (\emph{Ordinatio}~II, d.~6, q.~2) traces Lucifer's first sin to immoderate self-love (\emph{affectio commodi} / \emph{luxuria spiritualis}), whereas Francisco Su\'arez (\emph{De angelis}~VII.13) and later Scotists center his revolt on refusing the revealed Incarnation and hypostatic union (a.~3); St.~Augustine leaves first-instant voluntary sin undecided in \emph{City of God}~XI.15, while Aquinas's own mature \emph{De malo}~q.~16, a.~4 shifts the proof from efficient causality to cognitive non-simultaneity and excludes first-instant merit alongside sin (a.~5); St.~Bonaventure (\emph{Breviloquium}~II.7) and several Fathers allow a temporal interval (\emph{mora}) before the fall (a.~6); St.~John of Damascus (\emph{De fide orthodoxa}~II.4) holds that the fallen chief presided over the terrestrial realm rather than the highest seraphic rank (a.~7); none recorded on aa.~8--9.}

\subsection{The Punishment of the Fallen Angels (\ST{I, q.~64})}

Question~64 examines four dimensions of the demons' punishment: their intellect, their will, their sorrow, and their place. First, is the demonic intellect stripped of all knowledge of truth (\ST{I, q.~64, a.~1})? Pseudo-Dionysius rules out any such natural mutilation: ``certain gifts were bestowed upon the demons which, we say, have not been changed at all, but remain entire and most brilliant'' (\emph{De divinis nominibus}~IV). Aquinas distinguishes three kinds of truth-knowledge. \emph{Natural knowledge} is neither removed nor lessened in demons, because the simplicity of angelic substance admits no amputation of natural powers as a man might lose a hand or foot. \emph{Speculative knowledge of grace}---the revelation of divine secrets---is severely curtailed, disclosed to demons only selectively through holy angels or temporal workings of divine power (St.~Augustine, \emph{City of God}~IX.21; hence ``if they had known it, they would never have crucified the Lord of glory,'' 1~Cor~2:8; \ST{I, q.~64, a.~1, ad~4}). \emph{Affective knowledge of grace}---the gift of wisdom uniting the soul to God in love---is utterly lost. Demons still know creatures in ``nocturnal knowledge'' because they refuse to lift creation back to God's praise, combining keen natural insight, selective revelation, and long experience (\ST{I, q.~64, a.~1, ad~3, 5}); yet as \emph{De malo}~(q.~16, a.~6) shows in detail, their proud wills presumptuously judge above their light and plunge them into both speculative and practical error concerning supernatural mysteries.

Second, the will of the demons is irrevocably obstinate in evil (\ST{I, q.~64, a.~2}). Origen had speculated (\emph{Peri Archon}~I.6) that every created will remains perpetually mutable and that demons might one day be restored (\gk{apokatastasis}). The Church solemnly condemned that opinion at Constantinople in 543/553---anathematizing anyone who teaches that the punishment of demons is temporary and that an \gk{apokatastasis} will occur (Justinian's Edict, can.~9; Fifteen Anathematisms, can.~1)---and the Fourth Lateran Council confessed the ``perpetual punishment'' (\emph{poenam perpetuam}) awaiting the reprobate with the devil (DS~801). Scripture itself speaks of the demons' unyielding rebellion: ``the pride of them that hate thee ascendeth continually'' (Ps~73:23 [Heb~74:23]).

Aquinas seeks the cause of this obstinacy not in any limit of divine mercy---for mercy pardons the repentant, whereas demons cannot repent (\ST{I, q.~64, a.~2, ad~2})---but in the condition of angelic nature itself, citing St.~John of Damascus's maxim that ``death is to men, what the fall is to the angels'' (\emph{De fide orthodoxa}~II.3, where Damascene grounds angelic incapacity for repentance in incorporeality). Appetite always matches apprehension. Because human reason moves discursively from one angle to another, human choice remains flexible both before and after acting until death seals the soul; an angelic intellect, by contrast, grasps an object whole and immovably, as human intellect grasps first principles. Consequently, an angel's free-will is flexible before choice, yet clings immovably once it has chosen (\ST{I, q.~64, a.~2}; cf.~St.~Bonaventure, \emph{Breviloquium}~II.6--7). In \emph{De malo} (q.~16, a.~5), Aquinas refines this psychology through a threefold analysis of \emph{liberum arbitrium}---freedom among means to an end (which remains in demons), flexibility toward good or evil, and mutability of prior choice---showing that the force of their first choice endures in every later decision (\emph{permanet vis primae electionis ipsorum}). Notably, in \emph{De malo} (q.~16, a.~5, ad~13), he adds that by absolute power (\emph{secundum absolutam potentiam}) God could turn a demon's will back to good, yet doing so does not befit angelic nature. Every deliberate act of a demon is wicked---even telling truth to deceive or tremblingly confessing Christ under compulsion of evidence (Mark~1:24; Jas~2:19)---while their natural acts of existing and understanding remain naturally good, though twisted to malice (\ST{I, q.~64, a.~2, ad~5}).

Third, intense spiritual sorrow dwells in the demons (\ST{I, q.~64, a.~3}). Although pure spirits suffer no bodily passions, sorrow as a simple act of the will is the will's sharp resistance to what is or what is not (Apoc~18:7). Deprived of the happiness they naturally desire, curbed by God in their wicked designs, and consumed with envy over every soul saved, demons suffer constant volitional grief and tremble before coming judgment (Jas~2:19), even while devoid of reverential fear (\ST{I, q.~64, a.~3, ad~1--2}). Quoting St.~Augustine (\emph{City of God}~XIX.13), Aquinas notes a stark testimony to creation's goodness: ``sorrow for good lost by punishment, is the witness to a good nature''---demons never grieve for sin as sin, which would reveal a good will, only for punishment and lost happiness, involuntarily proving the goodness of the nature God gave them (\ST{I, q.~64, a.~3, ad~3}).

Fourth, until the Last Judgment, demons possess a twofold penal place: hell and this darksome atmosphere (\ST{I, q.~64, a.~4}). St.~Augustine observes that ``the darksome atmosphere is as a prison to the demons until the judgment day'' (\emph{De Genesi ad litteram}~III.10), reflecting St.~Paul's warning of ``the spirits of wickedness in the high places'' (Eph~6:12; cf.~2~Pet~2:4; St.~Bonaventure, \emph{Breviloquium}~II.7). Because Divine Providence procures human welfare both directly through holy angels and indirectly by allowing hostile spirits to test and exercise the faithful---preserving even fallen spirits within the utility of the cosmic order---some demons already torment the damned in hell while others dwell in the lower air until Judgment Day, after which all wicked angels and men will be confined to hell. A corporeal place punishes an incorporeal spirit not by changing its substance, but by afflicting its proud will through hateful confinement (\ST{I, q.~64, a.~4, ad~1}). Rejecting the error that punishment or glory is postponed until Judgment Day (2~Cor~5:1), Aquinas answers why roaming the air lightens none of their torment: just as a holy angel sent to earth loses no heavenly glory, like a bishop absent from his throne, demons carry the sentence of hell within themselves, for as the Gloss on James~3:6 declares, ``They carry fire of hell with them wherever they go'' (\ST{I, q.~64, a.~4, ad~3}). The Church's liturgical rites, canonical discipline, and pastoral ministry concerning demonic assault belong properly to the historical study of Catholic exorcism.

\angeldossier{\ST{I, q.~64, aa.~1--4} (Punishment of the demons: retention of natural knowledge with loss of grace, irrevocable obstinacy of will, volitional sorrow, and the twofold penal place until the Last Judgment).}{Common teaching (aa.~1, 3--4); Defined (a.~2).}{Gen~1:5; Job~41:24; Ps~73:23 [Heb~74:23]; Mark~1:24, 5:10; Luke~8:31; 1~Cor~2:8; 2~Cor~5:1, 6:14; Eph~5:13, 6:12; Jas~2:19, 3:6 (\emph{Glossa}); 2~Pet~2:4; Apoc~18:7; Constantinople (543/553) against Origenist \gk{apokatastasis}; Fourth Lateran Council, \emph{Firmiter credimus} (DS~801); Pseudo-Dionysius, \emph{De divinis nominibus}~IV; St.~Augustine, \emph{De Genesi ad litteram}~III.10, IV.22, V.19, \emph{De civitate Dei}~IX.21, XI.33, XIX.13, and \emph{De Genesi contra Manichaeos}~II.17; St.~Gregory the Great, \emph{Moralia in Iob}~IV; St.~John of Damascus, \emph{De fide orthodoxa}~II.3--4; St.~Bonaventure, \emph{Breviloquium}~II.7; St.~Thomas Aquinas, \emph{De malo}~q.~16, aa.~5--6.}{Origen's hypothesis of universal mutability and eventual restoration (\gk{apokatastasis}) of the demons, condemned at Constantinople in 543/553; in \emph{De malo}~q.~16, a.~5, ad~13, Aquinas notes that God could convert a demon's will by absolute power (\emph{potentia absoluta}), though this does not befit angelic nature (a.~2); none recorded on aa.~1, 3--4.}
